\documentclass[border=4pt]{standalone}
\usepackage{amsmath,amssymb,mathtools,xcolor,varwidth}
\definecolor{resultred}{HTML}{D81B60}
\begin{document}
\begin{varwidth}{28em}
\large\bfseries
As $P$ and $Q$ coincide, the ratio of $Qv$ to $Qx$ becomes a ratio of equality by      Cor. 2 of Lemma VII, so $Qx^2 \to 4PS \cdot QR$. From the similar triangles $QxT$ and      $SPN$, $Qx^2 : QT^2 = PS^2 : SN^2 = PS : SA$, whence \textcolor{resultred}{$QT^2 \propto SP \cdot QR$}.      By Cor. 1 of Prop. VI the force is measured as $1/(SP^2 \cdot QT^2/QR)$, and substituting      gives the \textcolor{resultred}{inverse-square law $F \propto 1/SP^2$}. The parabola obeys the same      law as the ellipse. \textit{Q.E.D.}
\end{varwidth}
\end{document}
